\begin{notechapter}
\chapter{Symmetric Operators and the SVD}
\label{ch:spectral-svd}
\chaptersubtitle{%
  Quadratic forms, eigenvalues and the shape of a linear map}

\begin{chapterplan}
\begin{itemize}[nosep]
  \item Self-adjoint operators (the notion depends on the inner product); the
        spectral theorem; Rayleigh quotients and the Courant--Fischer principle.
  \item Bilinear and quadratic forms as objects of their own; signature;
        positive definiteness.
  \item The singular value decomposition as the spectral theorem for $\adj AA$;
        operator norm, condition number, pseudo-inverse.
  \item Best low-rank approximation (Eckart--Young).
\end{itemize}
\smallskip\noindent\textbf{Machine learning threads.}
\begin{itemize}[nosep]
  \item Principal component analysis.
  \item Hessians, curvature, and why conditioning governs the speed of gradient
        descent.
  \item Low-rank adaptation of large models.
\end{itemize}
\end{chapterplan}
\end{notechapter}
